\documentclass{article}
\usepackage[a4paper, total={16cm, 26cm}]{geometry}
\usepackage{enumitem}
\usepackage{dashrule}

\def\trait{\vskip15pt\rule{\linewidth}{0.5pt}}

\begin{document}

\begin{center}
    \Large Évaluation de cours : Arithmétiques (Sujet A)
\end{center}
\vskip10pt
Nom : \hskip50pt Prénom : \hskip50pt Classe~:
\vskip20pt
\begin{enumerate}[label=\arabic*), font=\bfseries]
    \item Écrire la définition d'un nombre premier.
    \trait\trait
    \item Écrire tous les diviseurs de 20.
    \trait\trait
    \item Écrire trois multiples de 12.
    \trait
    \item Écrire la liste des nombre premiers inférieurs ou égaux à 30.
    \trait
    \trait
    \item Le nombre 0 est-il un nombre premier~? Justifier.
    \trait\trait\trait
    \item Décomposer les nombres suivants en facteurs premier~:
    \begin{enumerate}[label=\alph*), font=\bfseries]
        \item 30\vskip70pt
        \item 120\vskip70pt
        \item 63
    \end{enumerate}
\end{enumerate}
\newpage
\begin{center}
    \Large Évaluation de cours : Arithmétiques (Sujet B)
\end{center}
\vskip10pt
Nom : \hskip50pt Prénom : \hskip50pt Classe~:
\vskip20pt
\begin{enumerate}[label=\arabic*), font=\bfseries]
    \item Écrire la définition d'un nombre premier.
    \trait\trait
    \item Écrire tous les diviseurs de 45.
    \trait\trait
    \item Écrire trois multiples de 15.
    \trait
    \item Écrire la liste des nombre premiers inférieurs ou égaux à 30.
    \trait
    \trait
    \item Le nombre 0 est-il un nombre premier~? Justifier.
    \trait\trait\trait
    \item Décomposer les nombres suivants en facteurs premier~:
    \begin{enumerate}[label=\alph*), font=\bfseries]
        \item 12\vskip70pt
        \item 150\vskip70pt
        \item 72
    \end{enumerate}
\end{enumerate}
\end{document}